\subsection{The quaternion cover: the complete fixed subgroup}

\begin{statement}{A.36}{Proved}{All quaternion-cover fixed elements are in the selected cusp inertia}
For the quaternion first cover \(C_Q\setminus S\) of A.4,
every prime \(\ell\), every packet section induced as in
Part A, and every open subgroup \(\Delta_0\) preserving the
chosen cover and cusp lift, the fixed subgroup of the actual
geometric pro-\(\ell\) group is exactly the chosen cusp inertia:
\[
       \left(\pi_1((C_Q\setminus S)_{\bar K})^{(\ell)}
                          \right)^{s_b(\Delta_0)}
                              =I_{P,\ell}.              \tag{A.36.1}
\]
At odd primes the proof applies more generally to every
cusp packet on this punctured quartic over a finite number
field of definition, with its own derived Galois subgroup.
No nonzero Kummer condition on the packet is required.

The pointed stable reduction at \(2\) has two rational
components each carrying exactly two of the four cusps
and meeting the rest of the fibre at one node. For the
cyclic-eight map of A.35, the cusps above \(B=0,\infty\)
form one such pair; the two cusps above \(B=1\) form the
other. This pairing and the actual first ordinary
monodromy derivation, rather than injectivity of the
whole covering homomorphism on pro-\(\ell\) groups,
give the required comparison.
\end{statement}

\begin{proof}
\paragraph{The four-marked rational tree.}
We use the common-node specialization proved in A.35.
The original \(Q_8\) action preserves both components
adjacent to that node. On the metric edge of the unpointed
stable model, it fixes both ends, hence acts identically.
The four cusps have the same retraction point on this
edge, since they form a single \(Q_8\)-orbit. This point
is interior: the cusps specialize to the node, not to a
smooth point of either adjacent component.
After a finite extension, subdivide at that point.
The resulting rational bridge \(R\) has two attachment
points fixed by \(Q_8\).

Every element of the induced \(Q_8\) action on \(R\)
has order a power of two. A nonidentity automorphism
of \(\mathbf P^1\) of two-power order in characteristic
two is unipotent and has exactly one geometric fixed
point. It cannot fix both attachments. Thus the action
on \(R\) is trivial. Transitivity now implies that all
four marked sections meet the same smooth point of
\(R\) before their marked separation.
Their separation therefore attaches a rooted rational
tree to \(R\), with four marked points and no positive
genus component in this branch.

We record the short stability argument determining this
tree. Its root is fixed by \(Q_8\), so transitivity of
the four marks forces their marked vertices to have
the same distance from the root. A marked vertex
cannot have a descendant: such a descendant subtree
has a marked leaf, at strictly greater distance.
Hence all marked vertices are leaves. Every leaf has
at least two marks, by stability, while every unmarked
internal vertex has at least two outgoing edges.
With four marks, there are just two possibilities:
one component with four marks, or an unmarked root
component with two children, each carrying two marks.
There are no intervening unmarked vertices of valence
two in the stable tree.

The first possibility is excluded by the cyclic-eight
automorphism \(\tau\) from A.35. It preserves the root,
fixes the two cusps over \(B=0,\infty\), and exchanges
the other two. On a single four-mark component it would
fix the attaching point and two distinct marked points.
An automorphism of \(\mathbf P^1\) fixing three points
is the identity, contradicting that exchange. We
therefore obtain the two two-mark leaves.
The two \(\tau\)-fixed cusps lie on the same leaf:
if they lay on different leaves, \(\tau\) would preserve
both leaves and fix each leaf's other marked point as
well. This again contradicts the exchange.
Each marked leaf is consequently a three-punctured
rational component.

The metric statements used here have their elementary
local interpretation in \(uv=\varpi^e\): retraction
records \(v(u(P))\) on the interval \([0,e]\).
An automorphism preserving the two branches preserves
this coordinate, and equal coordinates produce the
same inserted bridge. Thus no choice of a complex
picture or of a conjectural reduction criterion
enters the marked-tree argument.

\paragraph{The first ordinary derivation and its packet ambiguity.}
Fix a cusp \(P\), and call the other cusp on its leaf
\(Q\); call the remaining pair \(R_1,R_2\).
Let \(r_1>0\) be the edge length from the \(P,Q\) leaf
to the unmarked root of this rational branch, and
\(r_2>0\) the length from that root to the other marked
leaf. All other edges lead toward the genus components.

The geometric pro-\(\ell\) group is free of rank nine.
Use six handle generators and three independent cusp
generators, with first homology classes
\[
         v_1,\ldots,v_6,\quad p,q,r.
\]
Here \(p,q,r\) correspond to \(P,Q,R_1\); the remaining
cusp has first class \(-p-q-r\). Put \(s=p+q\).
We use the \emph{ordinary} lower central degree in this
paragraph, so all nine displayed generators have degree
one. The first derivation of the edge twists, with a
base point on the \(P,Q\) leaf, is
\[
 \begin{split}
 D_0p&=D_0q=0,\\
 D_0v_j&=-r_1[s,v_j]\quad(1\le j\le6),\\
 D_0r&=-(r_1+r_2)[s,r].
 \end{split}                                             \tag{A.36.2}
\]
Indeed the first edge encloses all generators outside
the base leaf and its boundary has first class
\(-p-q=-s\). The edge toward the other marked leaf
has first boundary class \(r+(-p-q-r)=-s\).
Edges enclosing only genus components have boundary
in the commutator subgroup and contribute first in
ordinary degree two, not in (A.36.2). This verifies both
coefficients and their signs from the based edge-twist
formula of \cite[Definition 3.1.5 and Theorem 3.1.6]{fghBettsDogra}.

There is a small base-point issue worth making explicit:
that theorem is stated for ordinary points. Choose an
ordinary local point in the punctured residue disc of
\(P\), and blow up its cusp half-edge until it extends
as a smooth point. The theorem applies directly there.
The added segment along that half-edge contributes
only an inner multiple of \(\operatorname{ad}p\).
Comparison with a tangential fibre takes place inside
the punctured formal disc at \(P\). Its geometric
prime-to-two covers are the tame Kummer covers
\(z\mapsto z^{\ell^n}\); the comparison path's cocycle
therefore lies in the actual cusp inertia \(I_{P,\ell}\).
Changing to any completed tangent packet also changes
the action by an inner cusp cocycle. Thus the genuine
degree-one derivation has the form
\[
                   D_\lambda=D_0+\lambda\operatorname{ad}p,
                  \qquad \lambda\in\mathbf Q_\ell.       \tag{A.36.3}
\]
No value of \(\lambda\), and no nonvanishing assertion
about that value, is assumed.

Enlarge the global field to \(L\), so that the curve,
marked model, selected packet, and subsequent covering
map are defined, and so that
\(\Gamma=[G_L,G_L]\subset\Delta_0\).
The commutator of local Frobenius at two with a tame
inertia element having nonzero \(\ell\)-adic tame
parameter belongs to \(\Gamma\). Its first derivation
is a nonzero common multiple of (A.36.3).
Here \(\ell\ne2\); semistability makes the prime-to-two
inertia action tame, and the Frobenius--tame relation
multiplies the parameter by \(q_{\mathrm{res}}-1\ne0\).
After rescaling the derivation, (A.36.3) applies
with some \(\lambda\).

\paragraph{An unbounded free-Lie kernel.}
Inside the free associative algebra on \(p,s,v_1,\ldots,v_6,r\),
let \(B=\mathbf Q_\ell\langle p,s\rangle\).
The derivation on \(B\) is \(\lambda\operatorname{ad}p\).
On any of the seven other letters \(z\), it has the form
\[
                  D_\lambda z=[\lambda p+c_zs,z],
                 \qquad c_z\ne0.                        \tag{A.36.4}
\]
Group words by their ordered sequence of these seven
letters, leaving words in \(B\) in the intervening slots.
The component with an empty such sequence is \(B\).
For any other component, take its first outside letter
\(z\). On its first \(B\)-slot the operator increasing
that slot's degree is
\[
                      \lambda L_p+c_zR_s.                \tag{A.36.5}
\]
All later contributions preserve that first-slot degree.
If \(\lambda=0\), (A.36.5) is injective by right
multiplication by the free letter \(s\).
If \(\lambda\ne0\), count occurrences of \(p\):
the left summand raises their number while the right
summand does not. It is again injective.
The greatest first-slot degree in a finite expression
therefore shows that there is no kernel outside \(B\).
Within \(B\) the kernel is all of \(B\) for
\(\lambda=0\), and is \(\mathbf Q_\ell[p]\) otherwise.
Intersecting with the primitive free Lie algebra gives
\[
 \ker D_\lambda=
 \begin{cases}
    \mathbb L(p,q),&\lambda=0,\\
    \mathbf Q_\ell p,&\lambda\ne0 .
 \end{cases}                                             \tag{A.36.6}
\]
This proof applies in every degree and to arbitrary
nonzero values of the seven coefficients in (A.36.4).

\paragraph{The actual target map removes the remaining two-letter kernel.}
The \(Q_8\) deck action is transitive on the marked pairs.
Compose the cyclic-eight map of A.35 with a deck
automorphism so that \(P,Q\) map respectively to
\(0,\infty\). After enlarging \(L\) if needed this is
an actual morphism of punctured curves
\[
          f:C\setminus S\longrightarrow
                    \mathbf P^1\setminus\{0,1,\infty\}.
                                                        \tag{A.36.7}
\]
Both selected cusp ramification indices are eight.
On first homology, \(p,q\) map to \(8X,8Z\), where
\(X,Z\) are the classes at \(0,\infty\). They are a
basis of the target rational first homology.
Consequently the induced map on the free Lie algebra
\(\mathbb L(p,q)\) is injective in every degree.
This asserts injectivity only on that explicitly
identified leading algebra.

Let \(g\) be an actual element fixed by the specified
\(\Delta_0\). Its image under (A.36.7) is fixed by
\(\Gamma=\Delta_L\) in the induced cusp packet.
Statement A.34 puts this image in the actual inertia
at zero. Since eight is a unit in \(\mathbf Z_\ell\),
there is an actual integral power \(c_P^a\) of the
selected cusp generator with the same image.
Replace \(g\) by \(h=g c_P^{-a}\).
This element is still fixed by \(\Delta_0\) and maps
to one under \(f\).

The ordinary lower central filtration of a finite-rank
free pro-\(\ell\) group is separated and has torsion-free
free-Lie graded pieces over \(\mathbf Z_\ell\).
If \(h\ne1\), its first nonzero class survives rationally
and is killed by the genuine first derivation
\(D_\lambda\). By (A.36.6) that class lies in
\(\mathbb L(p,q)\), including when \(\lambda\ne0\).
But its image in the target graded Lie algebra must
vanish because \(f(h)=1\). This contradicts the
injectivity just proved. Thus \(h=1\) and
\(g\in I_{P,\ell}\). The opposite inclusion holds
because \(\Delta_0\) has cyclotomic character one.

This proves (A.36.1) for every odd prime, with no
integral Malcev lifting. At \(\ell=2\), A.22 applies
to the original \(Q_8\) Galois closure and gives the
same conclusion. The finite certificate
\texttt{code/a\_Q8\_cusp\_fixed.py} checks the degree-one
edge coefficients, the rank-two inertia map, and the
kernel dimensions in small degrees; the unbounded
argument is (A.36.5).
\end{proof}
